% TOP_PROBLEM: {"id": 1389, "title": "Eternal Domination vs Domination Number", "category_id": 3, "category": {"id": 3, "name": "graph_theory", "display_name": "Graph Theory", "description": "Problems involving graphs, networks, and their properties.", "slug": "graph-theory", "order_index": 3, "created_at": "2026-07-31T15:26:25.671Z"}, "status": "open"}
% FIRST_POSTED: null
% ENTRY_KIND: statement_only
% Imported from: batch06.tex; UnsolvedMath ID 1389
\documentclass[11pt]{article}
\usepackage[margin=25mm]{geometry}
\usepackage{fontspec}
\setmainfont{Latin Modern Roman}
\usepackage{amsmath,amssymb,mathtools}
\usepackage{unicode-math}
\setmathfont{Latin Modern Math}
\usepackage{xurl,hyperref}
\hypersetup{colorlinks=true,urlcolor=blue}
\setcounter{secnumdepth}{0}
\setlength{\parindent}{0pt}
\setlength{\parskip}{5pt}
\setlength{\emergencystretch}{3em}
\newcommand{\sref}[2]{\hyperlink{source:#1:#2}{[S#2]}}
\newcommand{\cataloguescope}[1]{\paragraph{Recorded target.}#1}
\begin{document}
% UnsolvedMath ID 1389; source catalogue code GRAPH-002
\setcounter{section}{1388}
\section{Eternal Domination vs Domination Number}\label{1389}
{\small\sffamily UnsolvedMath ID 1389 · Graph Theory}
\subsection{Definitions and mathematical statement}

\paragraph{Shared notation.}$\mathbb N=\{1,2,\ldots\}$, and $\mathbb Z,\mathbb Q,\mathbb R,\mathbb C$ denote the integers, rationals, reals and complex numbers. Any other conventions are specified locally.


Let $G=(V,E)$ be a finite simple undirected graph. A set $D\subseteq V$ dominates $G$ if every vertex lies in $D$ or is adjacent to a vertex of $D$; $\gamma(G)$ is the minimum size of a dominating set. In the one-guard-move eternal domination game, the defender initially occupies a dominating set with $k$ guards. At each turn the attacker chooses an unoccupied vertex $v$; exactly one guard moves along an edge to $v$, all other guards stay, and the new occupied set must dominate $G$. The least $k$ admitting a defense against every arbitrarily long sequence of attacks is $\gamma^\infty(G)$. Finally $\theta(G)$ is the minimum number of cliques partitioning $V$ (equivalently, covering $V$).
\paragraph{Statement recorded in the supplied source.}
Does there exist a finite simple graph $G$ such that
\[
 \gamma(G)=\gamma^\infty(G)<\theta(G)?
\]
\sref{1389}{2}
\subsection{Short English statement}
Can the ordinary domination number equal the one-guard-move eternal domination number while both are strictly smaller than the clique covering number?
\subsection{Sources}
\begin{description}
\item[\hypertarget{source:1389:1}{S1}] ULAM.AI and the UnsolvedMath Contributors, \emph{UnsolvedMath: A Curated Collection of Open Mathematics Problems}, record 1389 (GRAPH-002). CC BY 4.0. \url{https://huggingface.co/datasets/ulamai/UnsolvedMath}.
\item[\hypertarget{source:1389:2}{S2}] Tom Adamczewski and William F. Klostermeyer, \emph{A Counterexample to an Eternal Domination Conjecture} (2026), arXiv:2609.11500v1, Sections 1--3, especially Conjecture 3.1. \url{https://arxiv.org/abs/2609.11500v1}.
\end{description}
\label{end:1389}
\end{document}
